Quand on évapore les solutions aqueuses de certains composés ioniques, des
molécules d'eau peuvent rester insérées dans le réseau cristallin. Le composé
est alors qualifié d'hydraté. Le nombre de molécules d'eau est variable.

\textbf{Exemple~:}
Le sulfate de cuivre(II) pentahydraté~; sa formule s'écrit
\ce{CuSO4 * 5H2O}.
\begin{enumerate}
    \item Quelle masse de sulfate de cuivre(II) pentahydraté faut-il dissoudre
          dans l'eau pour obtenir $\qty{500}{\mL}$ de solution de concentration
          $\qty{0.040}{\mol\per\L}$~?
    \item Écrire l'équation de la dissolution de ce composé dans l'eau.
    \item Calculer la concentration des ions cuivre(II) et des ions sulfate dans
          cette solution.
\end{enumerate}

